\section{Introduction}
\label{sec:intro}

The field of deep learning has witnessed a significant paradigm shift toward decentralized and modular model training. In particular, \emph{model merging}—the practice of combining the weights of multiple independently fine-tuned neural networks without returning to the expensive pre-training stage—has emerged as a powerful mechanism to consolidate diverse task-specific capabilities \cite{modelsoups, taskarithmetic, tiesmerging}. Within the context of continual learning, model merging offers an elegant alternative to classic regularized or rehearsal-based strategies \cite{ewc, gem, der, lwf}, enabling practitioners to sequentially train experts on individual tasks and synthesize them post-hoc into a single multi-task model.

However, as model merging algorithms have advanced, they have grown increasingly complex. Modern frameworks often propose multi-stage, multi-component pipelines that claim "synergy" between specific optimization dynamics and post-hoc weight transformations. A representative example is \emph{Sharpness-Aware Isotropic Merging} (SAIM) \cite{saim2026}. SAIM argues that simple linear merging is insufficient for continual learning due to representation collapse in the shared parameter space. To resolve this, SAIM introduces a dual-stage solution: (1) a custom coordinate-wise sharpness-aware optimizer called Sharpness-Aware Block Coordinate Descent (SA-BCD) to guide fine-tuning, and (2) an SVD-based adaptive isotropic merging algorithm to balance the singular value spectrum of combined weight updates.

As \emph{methodologists}, we approach such multi-component claims with healthy scientific curiosity and skepticism. In machine learning, a recurring confounding variable is that complex frameworks are rarely evaluated against heavily tuned, strong baselines. When a multi-stage framework achieves superior performance, is it due to the complex, computationally expensive post-hoc transformations, or is it primarily an artifact of basic optimization hyperparameters (e.g., finding flatter minima) during individual expert training \cite{weight_averaging}? 

In this work, we systematically deconstruct the SAIM framework to isolate the individual contribution of each component and identify potential redundancies. We cross-evaluate five distinct optimization strategies and three merging strategies on a rigorous $5 \times 3$ grid, using Split CIFAR-100 \cite{cifar100} and a Vision Transformer backbone \cite{vit}. Our investigation yields several major, highly nuanced revelations:
\begin{enumerate}
  \item \textbf{Flatness is the Foundational Driver of Merging Success:} Under all evaluated mixing conditions, individual task-expert flatness is the primary driver of merging performance. Simply training experts with standard, globally perturbed Sharpness-Aware Minimization (SAM) \cite{sam} improves average accuracy by up to $+12.30\%$ over AdamW baselines.
  \item \textbf{The Benefit of SVD Merging is Highly Boundary-Condition-Sensitive:} Under the standard sequential fine-tuning parity regime ($\lambda = 0.0$), where no active parameter mixing occurs, we evaluate SVD as a boundary-condition sanity check. In this setting, SVD isotropic merging is mathematically redundant and acts as a distortive operator on un-mixed parameters, dropping accuracy by up to $12.7\%$. However, when active parameter-mixing is engaged ($\lambda = 0.2$), SVD-based Isotropic Merging acts as a highly effective regularizer against parameter interference, boosting SAM's performance to the overall best score of \textbf{76.42\%}.
  \item \textbf{Formulaic Bug and Suboptimality of SA-BCD:} We expose a fatal algebraic typo in SAIM's published SA-BCD optimizer formula that causes complete model divergence. Even when corrected, restricting sharpness perturbations to a selected coordinate subset (as done in SA-BCD) is empirically suboptimal compared to standard, globally perturbed SAM, which provides a $+5.4\%$ accuracy boost.
\end{enumerate}

By exposing these component-level dynamics, we show that while SVD-based post-hoc transformations are valuable for regularizing mixed parameter sets, they are entirely dependent on—and secondary to—optimizer-driven flatness. A simpler baseline—standard SAM training combined with naive Task Arithmetic merging—is highly competitive, and when SVD is employed, it should be paired with standard SAM rather than coordinate-restricted optimizers. Our work highlights the critical necessity of modular component-wise testing and rigorous boundary-condition analysis to establish a solid foundation for deep model merging.
